%% kapa: exam=1 round=22 year=2011 period=1 form=A subject=economics subject_ko=경제원론 seq=4 number=44
%% source: Q-Net 감정평가사 기출문제 (cst00309) artlSeq=5036790 file=extracted/1008/5036790/2011년 제22회 감정평가사 1차시험 기출문제/1교시.hwp
%% conversion: kordoc 4.24.51 -> tools/gen1.py (auto); review=proofread-sonnet; flags=none
\begin{question}{44}
저축의 역설(paradox of saving)에 관한 설명으로 옳은 것은?
\begin{choices}
  \item 소득이 증가하면 저축이 감소한다는 가설이다.
  \item 투자가 GDP와 정(＋)의 상관관계를 가질 때에는 저축이 증가하면 소득이 증가한다는 가설이다.
  \item 고전학파(Classical School)의 이론에서는 성립되지 않는 가설이다.
  \item 저축의 증가는 투자를 증가시킴으로써 경제성장을 촉진시킨다는 가설이다.
  \item 명목이자율의 상승이 인플레이션율을 하락시킨다는 가설이다.
\end{choices}
\end{question}
\answer{3}
